% Einheit 4 – Binomische Formeln (binomische-formeln.md Einheit 2) – Nr. 29–39
\einheitenkopf[e4][4 Binomische Formeln]{Einheit 4 von 5 · Binomische Formeln}
\zweigzeile{Hier lernst du, Quadrate von Klammern mit den binomischen Formeln auszumultiplizieren · neu in diesem Jahr · P10 · baut auf: Klammer mal Klammer (Einheit 3), Quadratzahlen}
\setcounter{aufgabe}{28}

\begin{aufgabe}{Ich erkenne, ob eine Klammer mit sich selbst malgenommen wird.}
\anweisung{Wird hier eine Klammer mit sich selbst malgenommen? Kreuze an.}
\begin{teilezwei}
\tz $(x + 6)^2$ \quad \janein & \tz $(x + 6) \cdot (x - 6)$ \quad \janein \\
\tz $(y - 2) \cdot (y - 2)$ \quad \janein & \tz $(z + 5) \cdot (5 + z)$ \quad \janein \\
\tz $(2x + 1) \cdot (x + 2)$ \quad \janein & \\
\end{teilezwei}
\end{aufgabe}

\begin{aufgabe}{Ich erkenne, welche binomische Formel passt.}
\anweisung{Kreuze an, welche binomische Formel passt. Rechne nichts aus.}
\begin{teile}
\teil $(x + 8)^2$ \hfill \kreuz{erste} \kreuz{zweite} \kreuz{dritte} \kreuz{keine}
\teil $(y - 9)^2$ \hfill \kreuz{erste} \kreuz{zweite} \kreuz{dritte} \kreuz{keine}
\teil $(x + 4) \cdot (x - 4)$ \hfill \kreuz{erste} \kreuz{zweite} \kreuz{dritte} \kreuz{keine}
\teil $(x + 2) \cdot (x + 7)$ \hfill \kreuz{erste} \kreuz{zweite} \kreuz{dritte} \kreuz{keine}
\teil $(z - 1) \cdot (z - 1)$ \hfill \kreuz{erste} \kreuz{zweite} \kreuz{dritte} \kreuz{keine}
\end{teile}
\end{aufgabe}

\begin{aufgabe}{Ich erkenne, was $a$ und was $b$ ist.}
\anweisung{In den Formeln steht $a$ für das erste und $b$ für das zweite Glied in der Klammer. Schreibe $a$ und $b$ auf. Rechne nichts aus.}
\begin{geruest}
\gz{a}{$(x + 9)^2$}{\feld{a} \quad \feld{b}}
\gz{b}{$(y - 4)^2$}{\feld{a} \quad \feld{b}}
\gz{c}{$(3x + 1)^2$}{\feld{a} \quad \feld{b}}
\gz{d}{$(5 + 2y)^2$}{\feld{a} \quad \feld{b}}
\gz{e}{$(4x - 7y)^2$}{\feld{a} \quad \feld{b}}
\end{geruest}
\end{aufgabe}

\verfahren{Ausmultiplizieren mit den Formeln}

\begin{aufgabe}{Ich kann mit den binomischen Formeln ausmultiplizieren.}
\anweisung{Multipliziere aus.}
\rechenplatz{2}
\begin{gleichungsraster}[2]
\gl{(x + 4)^2} & \gl{(x + 1)^2} \\
\gl{(y + 6)^2} & \gl{(x + 10)^2} \\
\gl{(x - 9)^2} & \gl{(x + 11) \cdot (x - 11)} \\
\gl{(x + 3) \cdot (x - 8)} & \\
\end{gleichungsraster}
\end{aufgabe}

\begin{aufgabe}{Ich kann mit den binomischen Formeln ausmultiplizieren – weiter.}
\begin{gleichungsraster}[2]
\gl{(3x + 4)^2} & \gl{(x - 5y)^2} \\
\gl{2 \cdot (x + 6)^2} & \gl{-(x - 8)^2} \\
\gl{(y - 4)^2 - 10} & \gl{(x - 6)^2 - 2 \cdot (3x - 4) + 7x \quad \text{(P10 2022 GYM)}} \\
\end{gleichungsraster}
\end{aufgabe}

\verfahren{Zwei Terme als gleich nachweisen}

\begin{aufgabe}{Ich kann nachweisen, dass zwei Terme gleich sind.}
\anweisung{Zeige durch Umformen, dass beide Terme gleich sind.}
\rechenplatz{2}
\begin{gleichungsraster}[2]
\gl{(x + 2)^2 - 1 \ \text{ und } \ x^2 + 4x + 3} & \gl{(x + 1)^2 + 7 \ \text{ und } \ x^2 + 2x + 8} \\
\gl{(x - 4)^2 + 2 \ \text{ und } \ x^2 - 8x + 18} & \\
\anweisung{Eine Parabel hat die Gleichung $y = (x + 6)^2 - 11$. Zeige, dass sie auch die Gleichung $y = x^2 + 12x + 25$ hat. (P10 2017 OS)}
\gl[3]{y = (x + 6)^2 - 11} & \\
\end{gleichungsraster}
\end{aufgabe}

\verfahren{Kopfrechnen mit den Formeln}

\begin{aufgabe}{Ich kann mit den binomischen Formeln im Kopf rechnen.}
\anweisung{Rechne mit einer binomischen Formel im Kopf.}
\begin{teilezwei}
\tz $21^2 = $ \leerfeld & \tz $51^2 = $ \leerfeld \\
\tz $19^2 = $ \leerfeld & \tz $98^2 = $ \leerfeld \\
\tz $32 \cdot 28 = $ \leerfeld & \tz $45 \cdot 55 = $ \leerfeld \\
\end{teilezwei}
\end{aufgabe}

\begin{aufgabe}{Ich finde den Fehler beim Ausmultiplizieren mit einer binomischen Formel.}
\anweisung{Finde den Fehler, benenne ihn und korrigiere die Rechnung.}
\begin{teile}
\teil Noah rechnet: \rechnung{(x + 6)^2 &= x^2 + 36}
\teil Ella rechnet: \rechnung{(3x - 2)^2 &= 3x^2 - 12x + 4}
\teil Finn rechnet: \rechnung{-(x + 9)^2 &= -x^2 + 18x + 81}
\end{teile}
\end{aufgabe}

\begin{aufgabe}{Ich kann entscheiden, welcher Term gleichwertig ist.}
\anweisung{Kreuze den Term an, der gleichwertig ist.}
\begin{teile}
\teil $(y + 9)^2$ \\
\kreuz{$y^2 + 81$} \\ \kreuz{$y^2 + 18y + 81$} \\ \kreuz{$y^2 + 9y + 81$} \\ \kreuz{$y^2 + 18y + 18$}
\teil $(x - 10) \cdot (x + 10)$ \\
\kreuz{$x^2 - 20x - 100$} \\ \kreuz{$x^2 - 100$} \\ \kreuz{$x^2 + 100$} \\ \kreuz{$x^2 - 20x + 100$}
\teil $(4 - 3y)^2$ \quad (P10 2019 GYM) \\
\kreuz{$16 - 9y^2$} \\ \kreuz{$16 - 24y + 9y^2$} \\ \kreuz{$16 - 12y + 9y^2$} \\ \kreuz{$16 - 24y - 9y^2$}
\end{teile}
\end{aufgabe}

\begin{aufgabe}{Ich kann begründen, warum beim Quadrat einer Summe ein Mittelglied entsteht.}
\begin{teile}
\teil Rechne $(6 + 2)^2$ und $6^2 + 2^2$. Begründe damit, warum man beim Quadrat einer Summe nicht jedes Glied einzeln quadrieren darf.
\teil Das Quadrat hat die Seitenlänge $x + 2$. Begründe am Bild, warum $(x + 2)^2$ nicht $x^2 + 4$ ist.
\end{teile}

\flaechenbild{x}{2}{x}{2}{}{}{}{}

\end{aufgabe}

\begin{aufgabe}{Ich kann eine Sachaufgabe mit einer binomischen Formel lösen.}
Ein quadratischer Spielplatz hat die Seitenlänge $x$\,m. Er wird nach rechts und nach unten um je $12$\,m vergrößert, sodass wieder ein Quadrat entsteht.

\flaechenbild{x}{12}{x}{12}{}{}{}{}

\begin{teile}
\teil Stelle einen Term für den Flächeninhalt des neuen Spielplatzes auf und multipliziere aus.
\teil Um wie viel m$^2$ wird der Spielplatz größer? Gib einen Term an.
\teil Der alte Spielplatz war $20$\,m lang. Wie viel m$^2$ kommen hinzu?
\end{teile}
\end{aufgabe}
